\chapter{Arbres et pricers aux différences finies en pratique}\label{ch:dv:trees-and-finite-difference-pricers-in-practice}

Un pricer renvoie 6.0864 pour un put américain à un an avec un arbre binomial à 200 pas, et 6.0975 avec 201 pas. Le trader
demande laquelle des deux valeurs il faut croire. La réponse est qu'aucune des deux n'est la bonne~: la valeur exacte, à cinq décimales près, est 6.0904, située sous l'une et au-dessus de l'autre. Quant
à l'\omterm{def:dv:the-binomial-model:crr}{arbre de Cox--Ross--Rubinstein} du chapitre 2, il converge vers elle, mais en zigzaguant selon la parité du nombre de pas,
et seulement au premier ordre. À 1\,500 pas, les deux voisins sont enfin à moins d'un centième de centime. Une grille de Crank--Nicolson avec le
démarrage adéquat et le strike sur un nœud y parvient avec 200 nœuds et 200 pas de temps, soit 28 fois moins de calcul. Ce chapitre porte sur
la différence entre un pricer qui converge et un pricer qui est précis. Il traite de l'équation de pricing sur une grille, des conditions
imposées par les dividendes et les barrières, de l'exercice américain par un algorithme direct, et de la manière de répartir un budget de calcul.

\section{L'équation de pricing sur une grille}

Pour un sous-jacent unique, l'\omterm{def:dv:black-scholes-three-ways:pde}{équation de Black--Scholes} (chapitre 3) se résout le plus naturellement en $x=\ln S$ et en temps restant jusqu'à l'échéance
$\tau$~:
\[
\partial_\tau V=\tfrac12\sigma^2\partial_{xx}V+\bigl(r-q-\tfrac12\sigma^2\bigr)\partial_xV-rV.
\]
Sur une grille $x_0<\dots<x_M$, les dérivées deviennent des différences à trois points. La construction autorise un pas variable, avec
$h_\pm$ les distances aux voisins, de sorte que les nœuds puissent être concentrés au voisinage du strike, là où le payoff présente son point anguleux.
En temps, le $\theta$-schéma (livre~4, chapitre 27) mélange des pas explicites et implicites. $\theta=\frac12$ correspond
au schéma de Crank--Nicolson, d'ordre deux et inconditionnellement stable, mais non exempt d'oscillations. Chaque pas résout un
système tridiagonal, en $O(M)$ opérations.

Crank--Nicolson présente une faiblesse connue, que les payoffs d'options déclenchent. Le schéma n'amortit que faiblement les composantes haute fréquence de la
solution. Un payoff à point anguleux contient toutes les fréquences, de sorte que les premiers pas laissent près du strike des
oscillations qui s'amortissent lentement. Les prix les remarquent à peine, mais les gammas, eux, oui. Le remède classique est le démarrage de Rannacher~: quelques demi-pas entièrement implicites
d'abord, qui amortissent les hautes fréquences, puis Crank--Nicolson. Un second remède agit sur le payoff lui-même.

\begin{definition}[Lissage du payoff]\label{def:dv:trees-and-finite-difference-pricers-in-practice:smoothing}
Le \emph{lissage du payoff}\index{lissage du payoff} remplace la valeur du payoff en chaque nœud de la grille par sa moyenne sur la cellule du
nœud, de sorte que la condition initiale porte l'effet du point anguleux sans sa dérivée discontinue~; combiné au démarrage de
Rannacher, il rétablit la convergence d'ordre deux du schéma pour les payoffs non réguliers.
\end{definition}

\begin{example}[Gamma au voisinage de l'échéance]\label{ex:dv:trees-and-finite-difference-pricers-in-practice:gamma}
Un put européen à une semaine de l'échéance (strike 100, volatilité de 20~\%), sur une grille de 200 nœuds avec dix pas de temps. Le
Crank--Nicolson nu donne une gamma de 0.217 au strike, où la valeur exacte est 0.141, et oscille autour de la courbe vraie de
part et d'autre. Avec deux pas de démarrage de Rannacher, l'erreur maximale est de 0.0002
(la \cref{fig:dv:trees-and-finite-difference-pricers-in-practice:gamma}). Un système de risque qui publie des gammas de Crank--Nicolson
sans ce démarrage montrera aux hedgers un pic de gamma fantôme à chaque strike proche de l'échéance.
\end{example}

\begin{omfigure}
\begin{tikzpicture}
\begin{axis}[width=10.5cm,height=5cm,
  xlabel={sous-jacent},ylabel={gamma},
  tick label style={font=\small},label style={font=\small},
  xmin=90,xmax=110,ymin=0,ymax=0.24,grid=major,grid style={gray!25},
  yticklabel style={/pgf/number format/fixed,/pgf/number format/precision=2},
  legend columns=3,legend cell align=left,legend style={font=\footnotesize,at={(0.5,-0.3)},anchor=north,draw=none,fill=none,/tikz/every even column/.append style={column sep=8pt}},
  table/col sep=comma]
\addplot[gray,very thick] table[x=s,y=exact]{figdata/derivatives/22-trees-and-finite-difference-pricers-in-practice/gamma.csv};
\addplot[omProp,thick,mark=*,mark size=1pt] table[x=s,y=cn]{figdata/derivatives/22-trees-and-finite-difference-pricers-in-practice/gamma.csv};
\addplot[omThm,very thick,dashed] table[x=s,y=rannacher]{figdata/derivatives/22-trees-and-finite-difference-pricers-in-practice/gamma.csv};
\legend{exacte,Crank--Nicolson,avec démarrage de Rannacher}
\end{axis}
\end{tikzpicture}

\omcaption{Gamma d'un put à une semaine sur une grille de 200 nœuds et dix pas de temps. Le Crank--Nicolson nu laisse une oscillation au
strike que dix pas ne peuvent amortir~; deux demi-pas implicites au départ la suppriment. Données~: le
tutoriel.}\label{fig:dv:trees-and-finite-difference-pricers-in-practice:gamma}
\end{omfigure}

Les arbres sont la même équation sur une grille particulière. L'arbre binomial lie le pas en espace au pas en temps, $\Delta x=\sigma\sqrt{\Delta t}$,
et constitue un schéma explicite. L'\omterm{def:dv:trees-and-finite-difference-pricers-in-practice:trinomial}{arbre trinomial} relâche ce lien.

\begin{definition}[Arbre trinomial]\label{def:dv:trees-and-finite-difference-pricers-in-practice:trinomial}
Un \emph{arbre trinomial}\index{arbre trinomial} déplace le log-prix vers le haut, vers le bas ou pas du tout à chaque pas, d'une quantité $\Delta x$ (par exemple
$\sigma\sqrt{3\Delta t}$), avec des probabilités choisies pour égaler la moyenne et la variance du log-rendement~; c'est un schéma aux différences finies explicite dont la
stabilité est intégrée à sa géométrie, et sa branche médiane supplémentaire permet de placer les nœuds de la grille plus librement qu'un
arbre binomial.
\end{definition}

\section{Dividendes discrets et conditions de saut}

Un dividende en numéraire $D$ versé en $t_d$ fait baisser l'action de $D$ à la date ex-dividende, tandis que la valeur de l'option est continue, puisque rien n'est
versé à son détenteur. La grille impose cela par une condition reliant les deux côtés de la date.

\begin{definition}[Condition de saut]\label{def:dv:trees-and-finite-difference-pricers-in-practice:jump}
Une \emph{condition de saut}\index{condition de saut} relie la solution d'une équation de pricing de part et d'autre d'une date à laquelle la variable d'état saute
de façon déterministe~: pour un dividende en numéraire, $V(t_d^-,S)=V(t_d^+,S-D)$, appliquée sur la grille par interpolation, suivie d'un nouveau démarrage de
Rannacher, car le profil interpolé cesse d'être régulier là où l'exercice américain intervient.
\end{definition}

Un arbre ne peut pas faire cela directement~: décaler les nœuds d'un \omterm{def:dv:the-binomial-model:recombining}{arbre recombinant} d'un montant en numéraire le fait cesser de se recombiner, ce pourquoi le chapitre 5 a utilisé le
modèle à dividende séquestré. La grille n'a pas ce problème.

\begin{example}[Un dividende de trois unités à six mois]\label{ex:dv:trees-and-finite-difference-pricers-in-practice:dividend}
Avec un dividende de 3 versé à six mois, le put européen à un an vaut 6.835 sur la grille, contre 6.719 par Black--Scholes appliqué au sous-jacent
diminué de la valeur actuelle du dividende, c'est-à-dire l'approximation à dividende séquestré à la même volatilité. Le put américain vaut 7.407, contre 6.090
sans dividende.
\end{example}

\section{Barrières et alignement sur la grille}

Une barrière surveillée en continu est une condition aux limites~: l'option vaut zéro, ou son remboursement, à la barrière. La grille doit donc
avoir un nœud exactement sur la barrière. Si la barrière tombe entre deux nœuds, le pricer la déplace soit vers le nœud le plus proche, ce qui introduit une
erreur d'ordre un sur le niveau de la barrière, soit interpole, ce qui ne fonctionne qu'approximativement. La grille étirée de la construction part de la
barrière et concentre les nœuds sur le strike. Les barrières surveillées de façon discrète s'appliquent comme une projection aux dates de surveillance (le
décalage du chapitre 15 en est la contrepartie analytique).

\begin{example}[Aligné et non aligné]\label{ex:dv:trees-and-finite-difference-pricers-in-practice:barrier}
Un call down-and-out à un an (strike 100, barrière 90) vaut 8.665 en forme fermée. Avec la barrière sur un nœud, l'erreur de la grille est de 0.0049
avec 40 nœuds et de 0.0002 avec 200. Avec des grilles uniformes sur lesquelles la barrière tombe entre deux nœuds, l'erreur est de 1.10 avec 40 nœuds, 0.16 avec 200 et
0.046 avec 300 (la \cref{fig:dv:trees-and-finite-difference-pricers-in-practice:barrier})~: des centaines de fois pire pour le même coût, et
dépendant de l'endroit précis où la barrière tombe entre deux nœuds.
\end{example}

\begin{omfigure}
\begin{tikzpicture}
\begin{axis}[width=10.5cm,height=5cm,xmode=log,ymode=log,
  xlabel={nœuds de la grille},ylabel={erreur absolue},
  tick label style={font=\small},label style={font=\small},
  xmin=35,xmax=340,ymin=5e-5,ymax=3,grid=major,grid style={gray!25},
  xtick={40,60,100,200,300},xticklabels={40,60,100,200,300},
  legend columns=2,legend cell align=left,legend style={font=\footnotesize,at={(0.5,-0.3)},anchor=north,draw=none,fill=none,/tikz/every even column/.append style={column sep=8pt}},
  table/col sep=comma]
\addplot[omThm,very thick,mark=*,mark size=1.6pt] table[x=m,y=aligned]{figdata/derivatives/22-trees-and-finite-difference-pricers-in-practice/barrier.csv};
\addplot[omProp,very thick,mark=triangle*,mark size=1.8pt] table[x=m,y=misaligned]{figdata/derivatives/22-trees-and-finite-difference-pricers-in-practice/barrier.csv};
\legend{barrière sur un nœud,barrière entre deux nœuds}
\end{axis}
\end{tikzpicture}

\omcaption{Erreur d'un call down-and-out à un an (prix d'exercice 100, barrière 90) sur des grilles de taille croissante, la barrière étant exactement sur un nœud et
entre deux nœuds. L'alignement, et non la taille, détermine la précision. Données~: le tutoriel.}\label{fig:dv:trees-and-finite-difference-pricers-in-practice:barrier}
\end{omfigure}

\section{Exercice américain}

L'exercice anticipé transforme l'équation de pricing en un problème à frontière libre (chapitre~6)~: $V\ge g$, la valeur d'exercice, partout, l'équation
étant vérifiée là où $V>g$. Sur une grille, le traitement naïf résout le système linéaire puis prend le maximum avec $g$. Cela fonctionne avec un pas
implicite, mais fait perdre à Crank--Nicolson sa précision, car la projection et la résolution ne sont pas cohérentes. Les méthodes itératives, telles que la sur-relaxation successive projetée,
résolvent correctement le système contraint, mais à un coût plusieurs fois supérieur. Pour un put, dont la région d'exercice est un unique intervalle
de cours bas, il existe une méthode directe.

\begin{definition}[Algorithme de Brennan--Schwartz]\label{def:dv:trees-and-finite-difference-pricers-in-practice:bs}
L'\emph{algorithme de Brennan--Schwartz}\index{algorithme de Brennan--Schwartz} résout un système tridiagonal sous la contrainte $V\ge g$ exactement
lorsque la région d'exercice est un intervalle à l'une des extrémités de la grille~: il élimine depuis l'autre extrémité, puis effectue la substitution
remontante vers la région d'exercice, en prenant le maximum avec $g$ à chaque nœud~; son coût est celui d'une seule résolution tridiagonale.
\end{definition}

L'ordre d'élimination constitue toute l'astuce. Pour un put, la région d'exercice se situe aux cours bas, si bien que l'élimination s'effectue depuis le haut
de la grille vers le bas, et la substitution remonte depuis la région d'exercice. Là, chaque valeur est fixée par la contrainte avant que les valeurs de continuation
situées au-dessus ne soient calculées. Pour un call sur une action versant un dividende, l'ordre s'inverse.

La grille fournit également la \omterm{def:dv:american-options-and-early-exercise:boundary}{frontière d'exercice} comme sous-produit~: le cours le plus élevé auquel la valeur est égale à la valeur d'exercice, à chaque date,
affiné entre les nœuds. L'arbre la fournit aussi, mais seulement sous la forme d'un escalier formé par ses propres nœuds
(la \cref{fig:dv:trees-and-finite-difference-pricers-in-practice:boundary}). Pour le put à un an, la grille place la frontière à 80.85 un an
avant l'échéance, alors que la frontière d'un arbre à 2\,000 pas oscille entre 80.68 et 81.04 d'un pas à l'autre.

\begin{omfigure}
\begin{tikzpicture}
\begin{axis}[width=10.5cm,height=5cm,
  xlabel={temps avant l'échéance (années)},ylabel={frontière d'exercice},
  tick label style={font=\small},label style={font=\small},
  xmin=0,xmax=1,ymin=78,ymax=101,grid=major,grid style={gray!25},
  legend columns=2,legend cell align=left,legend style={font=\footnotesize,at={(0.5,-0.3)},anchor=north,draw=none,fill=none,/tikz/every even column/.append style={column sep=8pt}},
  table/col sep=comma]
\addplot[omProp,thin] table[x=tau,y=s]{figdata/derivatives/22-trees-and-finite-difference-pricers-in-practice/boundary_tree.csv};
\addplot[omThm,very thick,mark=*,mark size=1.2pt] table[x=tau,y=s]{figdata/derivatives/22-trees-and-finite-difference-pricers-in-practice/boundary_grid.csv};
\legend{arbre binomial à 2\,000 pas,grille (400 nœuds)}
\end{axis}
\end{tikzpicture}

\omcaption{\omterm{def:dv:american-options-and-early-exercise:boundary}{Frontière d'exercice} du put américain (prix d'exercice 100, 5~\%, 20~\%) en fonction du temps restant avant l'échéance~: exercice en dessous de la courbe. La frontière de l'arbre est
un escalier de ses nœuds~; celle de la grille, affinée entre les nœuds, est lisse. Données~: le tutoriel.}\label{fig:dv:trees-and-finite-difference-pricers-in-practice:boundary}
\end{omfigure}

\section{Précision contre rapidité}

Tout pricer dispose de deux paramètres, la grille et le nombre de pas de temps, ainsi que d'un budget. La question que le trader posait dans l'introduction
revient en réalité à ceci~: combien de calcul faut-il pour obtenir une précision donnée~?

\begin{example}[Un dixième de cent]\label{ex:dv:trees-and-finite-difference-pricers-in-practice:budget}
Le put américain à un an (cours spot et prix d'exercice 100, taux 5~\%, volatilité 20~\%) vaut 6.0904, obtenu à partir d'une grille extrapolée par Richardson (1\,600 et 800 nœuds
et pas). L'erreur de l'\omterm{def:dv:the-binomial-model:crr}{arbre de Cox--Ross--Rubinstein} change de signe en alternance selon la parité du nombre de pas. Ses deux voisins sont tous deux à moins de 0.1 cent
à partir de 1\,500 pas, soit 1\,125\,750 mises à jour de nœuds. L'\omterm{def:dv:trees-and-finite-difference-pricers-in-practice:trinomial}{arbre trinomial} converge sans le zigzag, mais au premier ordre~: son erreur vaut 0.11 cent à 1\,600 pas
et 0.054 à 3\,200. La grille, avec un nœud sur le prix d'exercice, un \omterm{def:dv:trees-and-finite-difference-pricers-in-practice:smoothing}{lissage du payoff} et un démarrage de Rannacher, se trompe de 1.04 cent avec 50 nœuds et pas, de 0.30 avec
100, de 0.091 avec 200 et de 0.029 avec 400~: convergence du second ordre. Elle est à moins de 0.1 cent dès 200, soit 40\,600 mises à jour de nœuds. C'est 28 fois moins de travail que l'arbre
(\cref{fig:dv:trees-and-finite-difference-pricers-in-practice:work}). Sans lissage ni démarrage, la grille se trompe de 0.125 cent à 200.
\end{example}

\begin{omfigure}
\begin{tikzpicture}
\begin{groupplot}[group style={group size=2 by 1,horizontal sep=1.6cm},
  width=6.6cm,height=5cm,
  tick label style={font=\small},label style={font=\small},grid=major,grid style={gray!25},table/col sep=comma]
\nextgroupplot[xlabel={pas de l'arbre},ylabel={put américain},xmin=100,xmax=400,ymin=6.075,ymax=6.105,
  yticklabel style={/pgf/number format/fixed,/pgf/number format/precision=3},
  title={\small le zigzag binomial},
  legend columns=1,legend cell align=left,legend style={font=\footnotesize,at={(0.5,-0.32)},anchor=north,draw=none,fill=none}]
\addplot[omProp,thin] table[x=n,y=crr]{figdata/derivatives/22-trees-and-finite-difference-pricers-in-practice/oscillation.csv};
\addplot[black,thick,dashed] table[x=n,y=ref]{figdata/derivatives/22-trees-and-finite-difference-pricers-in-practice/oscillation.csv};
\legend{Cox--Ross--Rubinstein,valeur}
\nextgroupplot[xmode=log,ymode=log,xlabel={travail (mises à jour de nœuds)},ylabel={erreur (cents)},xmin=500,xmax=2e7,ymin=0.01,ymax=10,
  title={\small erreur en fonction du travail},
  legend columns=1,legend cell align=left,legend style={font=\footnotesize,at={(0.5,-0.32)},anchor=north,draw=none,fill=none}]
\addplot[omProp,thick,mark=*,mark size=1pt] table[x=work,y=err]{figdata/derivatives/22-trees-and-finite-difference-pricers-in-practice/work_crr.csv};
\addplot[omDef,thick,mark=square*,mark size=1.4pt] table[x=work,y=err]{figdata/derivatives/22-trees-and-finite-difference-pricers-in-practice/work_tri.csv};
\addplot[omThm,very thick,mark=triangle*,mark size=1.8pt] table[x=work,y=err]{figdata/derivatives/22-trees-and-finite-difference-pricers-in-practice/work_fd.csv};
\addplot[black,dashed] coordinates {(500,0.1) (2e7,0.1)};
\legend{binomial (le pire de $n$ et $n+1$),trinomial,grille (lissage et Rannacher),0.1 cent}
\end{groupplot}
\end{tikzpicture}

\omcaption{Put américain à un an. À gauche~: l'\omterm{def:dv:the-binomial-model:crr}{arbre de Cox--Ross--Rubinstein} comparé à la valeur, selon le nombre de pas. À droite~: erreur absolue, en cents,
en fonction du travail, en mises à jour de nœuds, pour l'arbre binomial (la plus grande erreur de chaque couple $n$, $n+1$), l'\omterm{def:dv:trees-and-finite-difference-pricers-in-practice:trinomial}{arbre trinomial} et
la grille de Crank--Nicolson. Données~: le tutoriel.}\label{fig:dv:trees-and-finite-difference-pricers-in-practice:work}
\end{omfigure}

Le travail, mesuré en mises à jour de nœuds, est un bon indicateur du temps d'exécution entre méthodes écrites dans le même langage, et le travail par nœud de la grille (une résolution tridiagonale
vaut quelques fois celui de l'arbre (une moyenne pondérée)~: l'accélération réelle est donc inférieure à 28, tout en restant du même ordre de grandeur. Deux autres économies
découlent de la convergence du second ordre. L'extrapolation de Richardson (livre~4) appliquée à deux tailles de grille élimine le terme d'erreur principal. L'étirement de la grille place les
nœuds là où le payoff présente son point anguleux. La construction intègre les deux. Les arbres conservent une place là où leur simplicité l'emporte sur leur coût~: l'enseignement,
la vérification, et les produits dont les règles d'exercice anticipé sont les plus faciles à énoncer nœud par nœud.

\section{Tutoriel~: un moteur aux différences finies}\label{tut:dv:trees-and-finite-difference-pricers-in-practice:tut}

\begin{tutorial}
\textbf{Objectif.} Construire un pricer de Crank--Nicolson sur une grille logarithmique du cours spot étirée, avec démarrage de Rannacher, \omterm{def:dv:trees-and-finite-difference-pricers-in-practice:smoothing}{lissage du payoff}, \omterm{def:dv:trees-and-finite-difference-pricers-in-practice:jump}{conditions de saut} liées aux dividendes,
barrières et exercice américain par Brennan--Schwartz~; le comparer aux arbres sur l'erreur en fonction du travail. \textbf{État final~:} les quatre
figures et les chiffres du problème du week-end.
\begin{tutsteps}
\item \textbf{L'opérateur sur une grille non uniforme}~: différences à trois points avec des pas inégaux~:
\omcode{code/firm/fdpricer/firm_fdpricer.py}{42}{50}{Coefficients tridiagonaux de l'opérateur de pricing.}\label{lst:dv:trees-and-finite-difference-pricers-in-practice:operator}
\item \textbf{Brennan--Schwartz}~: éliminer depuis le haut, puis substituer vers le haut sous la contrainte~:
\omcode{code/firm/fdpricer/firm_fdpricer.py}{68}{82}{Le solveur de Brennan--Schwartz.}\label{lst:dv:trees-and-finite-difference-pricers-in-practice:bs}
\item \textbf{La boucle temporelle}~: demi-pas de Rannacher, Crank--Nicolson, bords, dividendes~:
\omcode{code/firm/fdpricer/firm_fdpricer.py}{107}{135}{La boucle temporelle du theta-schéma avec conditions de saut.}\label{lst:dv:trees-and-finite-difference-pricers-in-practice:loop}
\item \textbf{Le jumeau C++20} du solveur, testé contre les valeurs Python à l'arrondi près (le jumeau Rust implémente le même algorithme)~:
\omcode{code/firm/fdpricer/cpp/firm_fdpricer.hpp}{36}{52}{Brennan--Schwartz en C++20.}\label{lst:dv:trees-and-finite-difference-pricers-in-practice:cpp}
\item \textbf{Exécuter} \texttt{dv\_fd.accuracy\_study()}, \texttt{gamma\_profile()}, \texttt{dividend\_example()}, \texttt{barrier\_alignment()} et
\texttt{fig\_fd.py}.
\end{tutsteps}
\textbf{Pour aller plus loin.} Remplacer Brennan--Schwartz par une résolution suivie d'une projection avec Crank--Nicolson et mesurer l'ordre de convergence~; pricer le
call américain avec dividende et repérer la \omterm{def:dv:american-options-and-early-exercise:boundary}{frontière d'exercice} de la grille juste avant la date ex-dividende~; chronométrer les versions Python, C++ et Rust sur la même
grille.
\end{tutorial}

\section{Construction~: le moteur aux différences finies}\label{bld:dv:trees-and-finite-difference-pricers-in-practice:fdpricer}

\begin{build}
\bfield{Rôle} Le moteur unidimensionnel d'EDP de la firme miniature~: le moteur aux différences finies de la bibliothèque de tarification du chapitre 28, utilisé pour
l'exercice américain, les barrières, les dividendes et les convertibles (chapitre 21)~; jumeaux en C++20 et en Rust pour la chaîne de production.
\bfield{Interface} \texttt{uniform\_grid(s0, vol, t, m, width)}~; \texttt{stretched\_grid(s0, k, vol, t, m, width, concentration, lower)}~;
\texttt{solve(payoff, x, t, r, q, vol, n\_t, american, rannacher, dividends, smoothing, lower, upper) -> values}~; \texttt{value\_at(x, v, s0) -> value,
delta, gamma}~; \texttt{american\_put}~; \texttt{smoothed\_payoff}~; \texttt{trinomial(s0, k, t, r, q, vol, n, right, american)}. C++20
\texttt{firm::fdpricer::american\_put}, Rust \texttt{firm\_fdpricer::american\_put}.
\bfield{Règles} Le strike (et toute barrière) est un nœud de la grille~; démarrage de Rannacher après l'échéance et après chaque discontinuité~; exercice américain par
l'\omterm{def:dv:trees-and-finite-difference-pricers-in-practice:bs}{algorithme de Brennan--Schwartz} lorsque la région d'exercice est un intervalle, par itération projetée sinon~; chaque prix est accompagné de son contrôle de convergence (deux
tailles de grille).
\bfield{Tests d'acceptation} \texttt{code/firm/fdpricer/tests/}~: valeur et delta européens contre Black--Scholes sur les deux grilles~; le put américain converge à
l'ordre deux~; l'\omterm{def:dv:trees-and-finite-difference-pricers-in-practice:trinomial}{arbre trinomial} converge~; la \omterm{def:dv:trees-and-finite-difference-pricers-in-practice:jump}{condition de saut} sur dividende contre une simulation~; une barrière placée sur un nœud contre la formule fermée~;
\texttt{cpp/} et \texttt{rust/}~: les jumeaux reproduisent les nombres du moteur Python à $10^{-9}$ près et convergent.
\bfield{Pour aller plus loin} Grilles bidimensionnelles (directions alternées) pour la volatilité stochastique~; repli par surrelaxation successive projetée~; pas de temps adaptatifs
autour des dividendes.
\end{build}

\begin{omsources}
\item M.~J.~Brennan et E.~S.~Schwartz, «~The valuation of American put options~», \emph{Journal of Finance} 32(2) (1977) 449--462.
\item R.~Rannacher, «~Finite element solution of diffusion problems with irregular data~», \emph{Numerische Mathematik} 43(2) (1984) 309--327.
\item D.~M.~Pooley, K.~R.~Vetzal et P.~A.~Forsyth, «~Convergence remedies for non-smooth payoffs in option pricing~», \emph{Journal of
Computational Finance} 6(4) (2003) 25--40.
\item D.~Tavella et C.~Randall, \emph{Pricing Financial Instruments: The Finite Difference Method} (Wiley).
\end{omsources}

\section{Exercices}

\begin{exercise}[$\star$]\label{exo:dv:trees-and-finite-difference-pricers-in-practice:1}
Pourquoi le prix de l'arbre binomial oscille-t-il au-dessus et au-dessous de la valeur lorsque le nombre de pas change de parité~?
\end{exercise}

\begin{exercise}[$\star$]\label{exo:dv:trees-and-finite-difference-pricers-in-practice:2}
L'erreur d'une grille vaut 0.30~cent avec 100 nœuds et 0.091 avec 200. Estimer l'ordre de convergence et l'erreur avec 400.
\end{exercise}

\begin{exercise}[$\star$]\label{exo:dv:trees-and-finite-difference-pricers-in-practice:3}
Pourquoi un \omterm{def:dv:the-binomial-model:recombining}{arbre recombinant} ne peut-il pas traiter directement un dividende en numéraire, et comment la grille le fait-elle~?
\end{exercise}

\begin{exercise}[$\star\star$]\label{exo:dv:trees-and-finite-difference-pricers-in-practice:4}
En utilisant les valeurs de la grille avec 800 et 1\,600 nœuds et pas, déduire la valeur extrapolée de Richardson pour une méthode d'ordre deux.
\end{exercise}

\begin{exercise}[$\star\star$]\label{exo:dv:trees-and-finite-difference-pricers-in-practice:5}
Expliquer pourquoi Brennan--Schwartz doit éliminer depuis le haut de la grille pour un put.
\end{exercise}

\begin{exercise}[$\star\star$]\label{exo:dv:trees-and-finite-difference-pricers-in-practice:6}
Pourquoi déplacer une barrière d'une demi-maille produit-il une erreur d'ordre un sur le prix~?
\end{exercise}

\begin{exercise}[$\star\star\star$]\label{exo:dv:trees-and-finite-difference-pricers-in-practice:7}
\emph{Programmation.} Pricer le put américain en résolvant par Crank--Nicolson puis en projetant, au lieu d'utiliser Brennan--Schwartz. Mesurer l'erreur avec 100, 200
et 400 nœuds, ainsi que l'ordre.
\end{exercise}

\begin{exercise}[$\star\star\star$]\label{exo:dv:trees-and-finite-difference-pricers-in-practice:8}
\emph{Trouver la faille.} «~Notre arbre price le put américain à six décimales~: nous l'avons lancé avec 2\,000 pas et le résultat n'a pas changé à la sixième décimale
d'une exécution à l'autre.~»
\end{exercise}

\section{Problème~: ni 200 ni 201}

\begin{problem}[{Problème du week-end --- combien de calcul coûte un dixième de cent}]\label{pb:dv:trees-and-finite-difference-pricers-in-practice:1}
Un système de risque price un portefeuille de puts américains à un an (spot et strike 100, taux 5~\%, volatilité 20~\%) avec un arbre binomial de 200 pas. Le desk doit certifier chaque
prix à un dixième de cent près.

\medskip\noindent\textbf{Partie I --- L'arbre.}
\begin{enumerate}
\item Donner le prix de l'arbre avec 200 puis 201 pas, ainsi que la valeur.
\item Donner les deux erreurs en cents.
\item À partir de combien de pas les deux voisins sont-ils tous deux à moins de 0.1~cent~?
\item Combien de mises à jour de nœuds cela représente-t-il~?
\item Pourquoi faire la moyenne de $n$ et $n+1$ n'est-il pas un vrai remède~?
\end{enumerate}

\medskip\noindent\textbf{Partie II --- La grille.}
\begin{enumerate}[resume]
\item Donner l'erreur de la grille avec 50, 100, 200 et 400 nœuds et pas.
\item Quelle taille atteint en premier la tolérance, et pour quel travail~?
\item Que perd la grille sans \omterm{def:dv:trees-and-finite-difference-pricers-in-practice:smoothing}{lissage du payoff} ni démarrage de Rannacher~?
\item Donner les erreurs de l'\omterm{def:dv:trees-and-finite-difference-pricers-in-practice:trinomial}{arbre trinomial} avec 1\,600 puis 3\,200 pas.
\item Que donnerait l'extrapolation de Richardson de 100 et 200 nœuds~?
\end{enumerate}

\medskip\noindent\textbf{Partie III --- Le reste du livre.}
\begin{enumerate}[resume]
\item Le livre contient aussi des puts avec un dividende de 3 dans six mois. Donner les prix européen et américain.
\item Et des calls down-and-out avec une barrière à 90~: qu'est-ce qui cloche sur une grille uniforme~?
\item Quelle grecque le Crank--Nicolson simple calcule-t-il mal près de l'échéance, et de combien dans l'exemple du chapitre~?
\item Pourquoi le travail en mises à jour de nœuds n'est-il qu'un indicateur du temps d'exécution~?
\item Comment le jumeau C++ modifierait-il la situation~?
\end{enumerate}

\medskip\noindent\textbf{Partie IV --- Jugement.}
\begin{enumerate}[resume]
\item Que faudrait-il dire au trader qui demande lequel de 200 et 201 croire~?
\item Comment on certifierait la précision d'un prix en production~?
\item Quand un arbre reste-t-il le bon outil~?
\item Énoncer le \emph{résultat nommé}~: la taille de grille et le travail pour lesquels le put américain est à moins de 0.1~cent, pour l'arbre et pour la grille, ainsi que l'accélération de
la grille.
\item En une phrase~: qu'est-ce qui distingue un pricer qui converge d'un pricer qui est précis~?
\end{enumerate}
\end{problem}

\section{Questions d'entretien}

\begin{interviewq}[$\star$ \iqroles{développeur, chercheur}]\label{iq:dv:trees-and-finite-difference-pricers-in-practice:1}
Pourquoi Crank--Nicolson produit-il des \omterm{def:dv:greeks-and-the-hedging-pnl:greeks}{grecques} oscillantes pour les options proches de l'échéance, et comment le corrigez-vous~?
\end{interviewq}

\begin{interviewq}[$\star\star$ \iqroles{développeur}]\label{iq:dv:trees-and-finite-difference-pricers-in-practice:2}
Comment pricez-vous un put américain sur une grille de différences finies de façon efficace~?
\end{interviewq}

\begin{interviewq}[$\star\star$ \iqroles{chercheur}]\label{iq:dv:trees-and-finite-difference-pricers-in-practice:3}
Comparez la convergence d'un arbre binomial, d'un \omterm{def:dv:trees-and-finite-difference-pricers-in-practice:trinomial}{arbre trinomial} et d'une grille de Crank--Nicolson pour un put américain.
\end{interviewq}

\begin{interviewq}[$\star\star$ \iqroles{développeur}]\label{iq:dv:trees-and-finite-difference-pricers-in-practice:4}
Comment traitez-vous un dividende en numéraire discret et une barrière surveillée en continu sur une grille~?
\end{interviewq}

\begin{interviewq}[$\star\star$ \iqroles{risque}]\label{iq:dv:trees-and-finite-difference-pricers-in-practice:5}
Comment valideriez-vous la précision numérique d'un pricer de production~?
\end{interviewq}

\begin{interviewq}[$\star\star\star$ \iqroles{développeur, chercheur}]\label{iq:dv:trees-and-finite-difference-pricers-in-practice:6}
Votre pricer en grille pour un call américain sur une action versant un dividende donne une \omterm{def:dv:american-options-and-early-exercise:boundary}{frontière d'exercice} légèrement erronée juste avant la date ex-dividende. Diagnostiquer.
\end{interviewq}
